\begin{proof}
Appendix~\ref{app:exact-norm} gives $F$, $\Delta$, $A_{\mathrm c}$, and
$B_{\mathrm c}$ in the present coordinates.  The seven coefficient
identities in Section~A.2 combine to give
\eqref{eq:norm-square-identity}.
The normalized residual curve is geometrically integral by
\cref{lem:normalized-function-field}, so $F$ is irreducible and $L$ is a
field.  The leading terms in A.2 give $\deg A_{\mathrm c}=3$ and
$\deg B_{\mathrm c}=2$ generically.  Reducing
\eqref{eq:norm-square-identity} modulo $F$ gives
\[
 -3\Delta=(A_{\mathrm c}/B_{\mathrm c})^2\ne0
\]
in $L$, proving \eqref{eq:cyclotomic-square-class} and showing that the
generic conic has rank exactly two.

For a rank-two symmetric $3\times3$ matrix $M$, one has
$\operatorname{adj}(M)=c_0aa^t$ for a nonzero vector $a$.  Hence the ratio
of any two nonzero principal cofactors is a square, so the square class of a
nonzero principal minor is intrinsic.  To see directly that this is the
ordering character, restrict the conic to the coordinate line $Y=0$.
If the two factors have coefficient rows $(a_0,a_1)$ and $(b_0,b_1)$ in
$(1,x)$, the resulting binary quadratic has discriminant
$(a_0b_1-a_1b_0)^2$; changing the sign of its square root interchanges the
two factors.  Thus the intrinsic quadratic algebra is obtained by
adjoining a square root of $\Delta$.  After shrinking to a normal open on
which $\Delta$ and $B_{\mathrm c}$ are units,
write $\varrho^2=-3$.  The maps
\[
 s\longmapsto\frac{A_{\mathrm c}}{\varrho B_{\mathrm c}},
 \qquad
 \varrho\longmapsto\frac{A_{\mathrm c}}{B_{\mathrm c}s}
\]
give inverse isomorphisms between the intrinsic algebra
$\mathcal O[s]/(s^2-\Delta)$ and the constant algebra
$\mathcal O[\varrho]/(\varrho^2+3)$.  This proves
\eqref{eq:cyclotomic-ordering-cover}.  The field extensions in
\eqref{eq:residual-function-fields} preserve the construction, and over
$\CC$ the constant quadratic cover is split.
\end{proof}

\begin{remark}[Arithmetic fibres]
The quadratic field $\QQ(\sqrt{-3})=\QQ(\zeta_3)$ explains the term
``cyclotomic.''  After spreading the rank-two open over a suitable
$\mathbf Z[1/N]$ with $6\mid N$, the two line factors at an
$\mathbf F_q$-point are
individually defined over $\mathbf F_q$ precisely when $-3$ is a square in
$\mathbf F_q$, equivalently when $q\equiv1\pmod3$.  If $q\equiv2\pmod3$,
they are exchanged by Frobenius and become defined over $\mathbf F_{q^2}$.
\end{remark}


\subsection{The scheme-theoretic Abel inverse}

Over the chart $T$ of \cref{lem:dominant-chart}, let
$p:C_T\to T$ be the relative curve and let $D$ be its universal divisor.
After restricting $T$, the sheaf
$p_*\mathcal O_{C_T}(3\kappa-D)$ is locally free of rank three.  The
residual equation $F_8=0$ defines a relative octic
\[
 \mathscr R_T\subset
 \PP_T\bigl(p_*\mathcal O_{C_T}(3\kappa-D)\bigr).
\]
After a further shrinking of $T$, the section $F_8$ is nonzero on every
projective-space fibre.  Thus $\mathscr R_T\to T$ is a flat projective
relative Cartier divisor.  Its generic fibre is geometrically integral by
\cref{lem:residual-octic-integral,lem:normalized-function-field}.  Flatness
forces every irreducible component to dominate the integral base $T$, so the
integral generic fibre gives a single component.  A nilpotent section would
vanish generically and hence be $\mathcal O_T$-torsion, which flatness
excludes.  Thus $\mathscr R_T$ is integral.

Let $\mathscr R_T^{(2)}\subset\mathscr R_T$ be the open subset on which the
residual conic has rank exactly two.  The relative Hilbert scheme of line
components of the universal conic is then finite \'etale of degree two over
$\mathscr R_T^{(2)}$.  Choosing one component determines the other and hence
orders them; denote this cover by
\[
 \operatorname{Ord}_{\mathrm{fac}}
    \longrightarrow\mathscr R_T^{(2)}.
\]
Its points will be written $([s],\ell_1,\ell_2)$, with
$Q_s=\ell_1\ell_2$ up to a unit pulled back from the base.

\begin{definition}[Admissible residual-norm locus]
\label[definition]{def:admissible-locus}
Let $\mathscr G\subset\mathscr R_T^{(2)}$ be the open locus on which the
following conditions hold after pullback to
$\operatorname{Ord}_{\mathrm{fac}}$:
\begin{enumerate}[label=\textup{(\roman*)}]
\item the fixed divisor $D$ is reduced;
\item the fixed norm line and the two lines $\ell_i=0$ meet $E$ transversely,
      their degree-three line sections are pairwise disjoint, and the residual
      line sections avoid the branch divisor of $C\to E$;
\item above each residual point, $s$ vanishes on exactly one sheet of
      $C\to E$, and every such zero is simple;
\item the two resulting degree-three zero divisors have $h^0=1$, and their
      Abel classes lie in the relative smooth locus of $\cX$.
\end{enumerate}
These conditions define an open subset of
$\operatorname{Ord}_{\mathrm{fac}}$ invariant under interchanging
$\ell_1$ and $\ell_2$.  Its image is open because the finite \'etale cover is
open, and invariance makes it the full inverse image of that image.  This
defines $\mathscr G$.
\end{definition}

\begin{lemma}[Relative ordered-zero construction]
\label[lemma]{lem:scheme-abel-inverse}
Put
\[
 \operatorname{Ord}_{\mathrm{fac},\mathscr G}
 =\operatorname{Ord}_{\mathrm{fac}}
   \times_{\mathscr R_T^{(2)}}\mathscr G.
\]
Let
$\operatorname{Ord}_{\mathrm{div}}$ be the functor over $\mathscr G$ whose
$W$-points are ordered pairs of relative effective Cartier divisors
$(D_1,D_2)$ of degree three on $C_W/W$ such that
\[
 \operatorname{div}(s_W)=D_W+D_1+D_2
\]
and the unordered pair $\{\pi_*D_1,\pi_*D_2\}$ is the unordered pair of line
sections of $E_W$ defined by the two components of $Q_{s_W}$.  This functor is
represented by a finite \'etale double cover of $\mathscr G$.
The assignment $Q_s=\ell_1\ell_2\mapsto(D_1,D_2)$ is an isomorphism of finite
\'etale double covers
\[
 \operatorname{Ord}_{\mathrm{fac},\mathscr G}
 \xrightarrow{\sim}\operatorname{Ord}_{\mathrm{div}}.
\]
It commutes with arbitrary base change, including nonreduced base change.
Moreover, each $D_i$ is a relative effective Cartier divisor of degree three,
and $[\mathcal O(D_i)\otimes\kappa^{-1}]$ lies in $\cX_{\mathrm{sm}/\cB}$.
\end{lemma}

\begin{proof}
For the forward construction put
$T'=\operatorname{Ord}_{\mathrm{fac},\mathscr G}$ and write
$Q_s=\ell_1\ell_2$ in its universal ordering.  By definition of
$\mathscr G$, none of the tangency, branch, collision, multiple-zero, or Abel
speciality loci excluded in \cref{def:admissible-locus} occurs.  For the
remainder of the forward construction, every object with a subscript $T$ is
pulled back to $T'$; this base change is suppressed from the notation.
The line $\ell_i$ then
cuts out a relative finite \'etale divisor $B_i\subset E_T$ of degree three.
Moreover
\[
 U_i=C_T\times_{E_T}B_i\longrightarrow B_i
\]
is finite \'etale of degree three, because $B_i$ avoids the branch divisor.

